% ═══════════════════════════════════════════════════════════════════════════════
% CHAPTER 6: SETUP AND DEPLOYMENT
% ═══════════════════════════════════════════════════════════════════════════════
\chapter{Setup and Deployment}

\section{Production: One-Shot Droplet Provisioning}

Production runs on a single DigitalOcean droplet provisioned by the helper script \texttt{scripts/}\linebreak[1]\texttt{provision-droplet.sh}. From a fresh Ubuntu 22.04 root shell:

\begin{infobox}[title=Provisioning Command]
\begin{verbatim}
DROPLET_IP="<DROPLET_IP>"
SITE_HOST="${DROPLET_IP//./-}.nip.io"

ssh -i ~/.ssh/tahkik_do root@$DROPLET_IP \
  "curl -fsSL https://raw.githubusercontent.com/\
mohamedbenhadjer/tahkik-server/main/scripts/provision-droplet.sh \
  | SITE_HOST=$SITE_HOST bash"
\end{verbatim}
\end{infobox}

The script is idempotent and does:

\begin{enumerate}
  \item Create a 4\,GB swapfile and tune \texttt{vm.swappiness=10}.
  \item Install Docker Engine + Compose plugin (with JSON log rotation).
  \item Detect a GPU (via \texttt{lspci}) and install the NVIDIA Container Toolkit if present.
  \item Clone the repository to \texttt{/opt/tahkik-server}.
  \item Write \texttt{.env} with \texttt{SITE\_HOST}, \texttt{MODEL\_ID}, \texttt{NEMO\_DEVICE}, \texttt{HEALTH\_TOKEN} (auto-generated if not supplied), and optional \texttt{HF\_TOKEN}.
  \item Build all three containers and bring the stack up. On a GPU droplet the \texttt{docker-compose.gpu.yml} overlay is applied automatically.
\end{enumerate}

First boot takes 10--15 minutes: \texttt{nemo\_toolkit[asr]} pulls roughly 1.5\,GB of wheels and the model checkpoint is around 500\,MB.

\section{Continuous Deployment}

The repository's \texttt{.github/workflows/deploy.yml} watches \texttt{main} and triggers a redeploy via SSH using the \texttt{appleboy/ssh-action} runner. On the droplet it runs:

\begin{lstlisting}[caption={GitHub Actions redeploy script}]
cd /opt/tahkik-server
git fetch origin main
git reset --hard origin/main
export GIT_SHA=$(git rev-parse --short HEAD)
export DEPLOYED_AT=$(date -u +%Y-%m-%dT%H:%M:%SZ)
docker compose pull caddy || true
if lspci | grep -i nvidia >/dev/null 2>&1; then
  docker compose -f docker-compose.yml -f docker-compose.gpu.yml \
                 up -d --build
else
  docker compose up -d --build
fi
docker image prune -f
\end{lstlisting}

\texttt{GIT\_SHA} and \texttt{DEPLOYED\_AT} are baked into the Go binary via \texttt{-ldflags} so the \texttt{/api/v1/system-health} response advertises the running build.

\section{Local Development}

For local iteration the Go and Python servers can run without Docker:

\phaseblock{Step 1 --- Python Environment}{%
\texttt{cd python}\\
\texttt{python3 -m venv .venv \&\& source .venv/bin/activate}\\
\texttt{pip install -r requirements.txt}
}

\phaseblock{Step 2 --- Start Inference Server}{%
\texttt{python inference\_server.py --port 8081}\\
Wait for the \texttt{[inference] READY} line on stdout.
}

\phaseblock{Step 3 --- Start Go Server}{%
\texttt{export INFERENCE\_SERVER\_URL=http://localhost:8081}\\
\texttt{export PORT=8080}\\
\texttt{go run ./cmd/api}
}

The local Python entry point (\texttt{python/inference\_server.py}) is the legacy Whisper development server using \texttt{benhadjermed/tahkik-small-warsh}. The production stack always uses the NeMo container at \texttt{python/hf\_space/}.

\section{Environment Variables}

\begin{table}[H]
  \centering
  \caption{Configuration environment variables}
  \label{tab:env_vars}
  \renewcommand{\arraystretch}{1.3}
  \begin{tabular}{llp{5.5cm}}
    \toprule
    \textbf{Variable} & \textbf{Default} & \textbf{Purpose} \\
    \midrule
    \texttt{SITE\_HOST}             & ---                                 & Public hostname Caddy issues a TLS cert for \\
    \texttt{DASHBOARD\_ORIGIN}      & \texttt{http://localhost:5173}      & CORS origin allowed by Caddy \\
    \texttt{MODEL\_ID}              & \texttt{nvidia/stt\_ar\_fc\dots}    & NeMo model identifier \\
    \texttt{NEMO\_DEVICE}           & \texttt{cpu}                        & \texttt{cpu} or \texttt{cuda} \\
    \texttt{TORCH\_NUM\_THREADS}    & \texttt{4}                          & CPU thread cap (irrelevant on GPU) \\
    \texttt{INFERENCE\_SERVER\_URL} & \texttt{http://localhost:8081}      & Backend URL seen by the Go server \\
    \texttt{HF\_TOKEN}              & ---                                 & Bearer token for HF access (private repos) \\
    \texttt{HEALTH\_TOKEN}          & ---                                 & Bearer token for \texttt{/api/v1/system-health} \\
    \texttt{PORT}                   & \texttt{8080}                       & Go server listen port \\
    \texttt{COMPOSE\_PROJECT\_NAME} & \texttt{tahkik-server}              & Compose label used by container enumeration \\
    \bottomrule
  \end{tabular}
\end{table}

\section{Hardware Recommendations}

\begin{table}[H]
  \centering
  \caption{Tested hardware tiers}
  \label{tab:hardware}
  \renewcommand{\arraystretch}{1.3}
  \begin{tabular}{llcc}
    \toprule
    \textbf{Tier} & \textbf{Hardware} & \textbf{Cold Start} & \textbf{Partial / Final p50} \\
    \midrule
    Development    & CPU laptop              & ${\sim}30$\,s & ${\sim}5$ / ${\sim}10$\,s \\
    Production CPU & DO \texttt{s-4vcpu-8gb} & ${\sim}60$\,s & 3--8 / 8--15\,s \\
    Production GPU & NVIDIA T4 / L4          & ${\sim}40$\,s & ${<}1$ / 1--2\,s \\
    \bottomrule
  \end{tabular}
\end{table}

% ═══════════════════════════════════════════════════════════════════════════════
% CHAPTER 7: API REFERENCE
% ═══════════════════════════════════════════════════════════════════════════════
\chapter{API Reference}

\section{Go Server Endpoints (Client-Facing)}

\subsection{POST /api/v1/evaluate}

Upload an audio file for Arabic Quranic recitation transcription, optionally with a reference text to enable madd analysis.

\begin{infobox}[title=Request Format]
Content-Type: \texttt{multipart/form-data}\\[2pt]
Required field: \texttt{audio} (file --- \texttt{.wav}, \texttt{.mp3}, \texttt{.m4a}, \texttt{.flac}, \texttt{.ogg}; max 50\,MB)\\[2pt]
Optional field: \texttt{reference\_text} (string --- bare text \emph{or} JSON-encoded array \texttt{[\{"surah":1,"ayah":1,"text":"\dots"\}, \dots]})
\end{infobox}

\begin{emeraldbox}[title=Success Response (200 OK) --- without reference\_text]
\begin{verbatim}
{
  "status": "success",
  "data": {
    "transcription": "<Arabic transcription text>",
    "confidence_score": 0.9423,
    "processing_time_ms": 1350
  }
}
\end{verbatim}
\end{emeraldbox}

\begin{emeraldbox}[title=Success Response (200 OK) --- with reference\_text]
\begin{verbatim}
{
  "status": "success",
  "data": {
    "transcription": "...",
    "confidence_score": 0.93,
    "processing_time_ms": 2100,
    "haraka_ms": 280,
    "errors": [
      {
        "word": "...",
        "word_index": 3,
        "rule": "Al-Madd",
        "sub_rule": "...",
        "severity": "minor",
        "expected": "2 counts",
        "detected": "~1.4 counts",
        "explanation": "Natural elongation of 2 counts...",
        "timestamp": {"start": 1.872, "end": 2.144}
      }
    ]
  }
}
\end{verbatim}
\end{emeraldbox}

\begin{warnbox}[title=Error Response (4xx / 5xx)]
\begin{verbatim}
{ "status": "error", "message": "unsupported audio format: .xyz" }
\end{verbatim}
\end{warnbox}

\subsection{GET /api/v1/system-health (Bearer-token gated)}

Returns the aggregated host + container + inference snapshot described in Section~\ref{subsec:system_health}.

\begin{emeraldbox}[title=Response (200 OK) --- shape only]
\begin{verbatim}
{
  "status": "ok",                       // ok | loading | degraded
  "host": {
    "hostname": "tahkik-nemo",
    "uptime_seconds": 81234,
    "cpu_percent": 14.2,
    "mem_used_mb": 3210,  "mem_total_mb": 7890,
    "swap_used_mb": 0,    "swap_total_mb": 4096,
    "disk_used_gb": 27,   "disk_total_gb": 160,
    "load_1m": 0.71, "load_5m": 0.65
  },
  "containers": [
    { "name": "caddy",     "status": "running",  "up_since": "..." },
    { "name": "go-server", "status": "healthy",  "up_since": "..." },
    { "name": "python",    "status": "healthy",  "up_since": "..." }
  ],
  "inference": { "requests_total": ..., "partial_p95_ms": ..., ... },
  "model":     { "id": "...", "device": "cpu", "loaded": true, ... },
  "version":   { "git_sha": "abc1234", "deployed_at": "2026-..." }
}
\end{verbatim}
\end{emeraldbox}

\subsection{GET /api/v1/system-health/ws (WebSocket)}

WebSocket upgrade of the above endpoint. Authentication is via \texttt{?token=<HEALTH\_TOKEN>} on the URL because browsers cannot set custom headers on a WebSocket handshake. The server pushes a fresh snapshot every 2 seconds and responds to client close frames cleanly.

\subsection{GET /health}

\begin{emeraldbox}[title=Health Check Response (200 OK)]
\begin{verbatim}
{"status": "ok"}
\end{verbatim}
\end{emeraldbox}

\section{Response Fields}

\begin{table}[H]
  \centering
  \caption{\texttt{data} object fields on \texttt{/api/v1/evaluate}}
  \label{tab:api_fields}
  \renewcommand{\arraystretch}{1.3}
  \begin{tabular}{llp{6.8cm}}
    \toprule
    \textbf{Field} & \textbf{Type} & \textbf{Description} \\
    \midrule
    \texttt{transcription}      & string & Arabic transcription of the recitation \\
    \texttt{confidence\_score}  & float  & Mapped RNNT log-prob score in $[0, 1]$ \\
    \texttt{processing\_time\_ms} & int  & Time in model (ms), excluding network \\
    \texttt{haraka\_ms}         & int    & Auto-calibrated one-haraka length (only when \texttt{reference\_text} was sent) \\
    \texttt{errors}             & list   & Detected \texttt{TajweedError} objects (only when \texttt{reference\_text} was sent) \\
    \bottomrule
  \end{tabular}
\end{table}

\section{Flutter Integration Example}

\begin{lstlisting}[language=Java, caption={Dart/Flutter API call example}]
final dio = Dio(BaseOptions(baseUrl: 'https://<site-host>'));

Future<Map<String, dynamic>> evaluate(File audioFile,
                                      {String? referenceText}) async {
  final formData = FormData.fromMap({
    'audio': await MultipartFile.fromFile(
      audioFile.path,
      filename: path.basename(audioFile.path),
    ),
    if (referenceText != null) 'reference_text': referenceText,
  });
  final response = await dio.post('/api/v1/evaluate', data: formData);
  return response.data as Map<String, dynamic>;
}
\end{lstlisting}

% ═══════════════════════════════════════════════════════════════════════════════
% CHAPTER 8: TROUBLESHOOTING
% ═══════════════════════════════════════════════════════════════════════════════
\chapter{Troubleshooting}

\subsection{Inference Server Unreachable / 502 from Caddy}

\phaseblock{Diagnosis}{Caddy is up but the Go server cannot reach the Python container. Either the Python container is still loading the NeMo model (cold start can take 10\,min on first boot --- the wheel install plus model download) or the container has crashed.}

\phaseblock{Resolution}{\texttt{docker compose -f /opt/tahkik-server/docker-compose.yml logs python --tail=50}. Watch for \texttt{[inference] model ready (warmup <ms>)}. If you see Python-level tracebacks, restart with \texttt{docker compose restart python}.}

\subsection{Model Download Timeout}

\phaseblock{Diagnosis}{First-time model download from HuggingFace can take several minutes on a slow link. The Python server prints \texttt{[inference] loading NeMo model \dots} and stays silent until the checkpoint is fully fetched into \texttt{/cache/nemo}.}

\phaseblock{Resolution}{Wait for the download to complete. Subsequent restarts use the cached weights in the \texttt{nemo\_cache} named volume.}

\subsection{system-health Returns 503 / 401}

\phaseblock{Diagnosis}{\texttt{503} means \texttt{HEALTH\_TOKEN} is unset in the environment (the endpoint fails closed). \texttt{401} means the \texttt{Authorization: Bearer <token>} header doesn't match the value baked into \texttt{.env}.}

\phaseblock{Resolution}{Verify \texttt{grep HEALTH\_TOKEN /opt/tahkik-server/.env} on the droplet and ensure the same value is wired into the dashboard's \texttt{VITE\_TAHKIK\_HEALTH\_TOKEN}.}

\subsection{Partial Latency Above the Window}

\phaseblock{Diagnosis}{Each partial transcribes 8\,s of audio. If inference itself takes longer than 8\,s, the stream falls behind real-time and the user experience degrades.}

\phaseblock{Resolution}{In order of preference: (1) reduce \texttt{PARTIAL\_WINDOW\_S} in \texttt{python/hf\_space/main.py} to 5; (2) raise \texttt{TORCH\_NUM\_THREADS} to the droplet's vCPU count; (3) move to a GPU droplet and apply \texttt{docker-compose.gpu.yml}.}

% ═══════════════════════════════════════════════════════════════════════════════
% CHAPTER 9: PERFORMANCE METRICS
% ═══════════════════════════════════════════════════════════════════════════════
\chapter{Performance Metrics}

\section{Inference Latency by Hardware}

The inference latency varies significantly between CPU and GPU droplets and between the partial-window and final inference paths. Measurements below were taken on the production droplet (\texttt{s-4vcpu-8gb}) for the CPU rows and on a single NVIDIA T4 for the GPU rows, using a 10-second Quranic audio sample (16\,kHz mono 16-bit WAV).

\begin{table}[H]
  \centering
  \caption{Average inference latency by hardware (p50 / p95)}
  \label{tab:perf_latency}
  \renewcommand{\arraystretch}{1.3}
  \begin{tabular}{llccc}
    \toprule
    \textbf{Path} & \textbf{Hardware} & \textbf{p50} & \textbf{p95} & \textbf{Throughput} \\
    \midrule
    Partial (8\,s window) & CPU 4-vcpu    & ${\sim}3.5$\,s & ${\sim}7.0$\,s & ${\sim}0.28$\,req/s \\
    Final (10\,s audio)   & CPU 4-vcpu    & ${\sim}8.5$\,s & ${\sim}15$\,s  & ${\sim}0.12$\,req/s \\
    Partial (8\,s window) & T4 GPU (FP16) & ${\sim}0.6$\,s & ${\sim}0.9$\,s & ${\sim}1.7$\,req/s  \\
    Final (10\,s audio)   & T4 GPU (FP16) & ${\sim}1.2$\,s & ${\sim}1.8$\,s & ${\sim}0.85$\,req/s \\
    \bottomrule
  \end{tabular}
\end{table}

\begin{infobox}[title=Key Observations]
\begin{itemize}
  \item The RNNT decoder dominates wall-time on CPU; GPU acceleration yields a ${\sim}6$--$8\times$ speedup on the same checkpoint.
  \item Final inference is ${\sim}2$--$3\times$ slower than the partial because it processes the full (capped at 60\,s) buffer rather than the trailing 8\,s.
  \item The Go API gateway introduces negligible overhead ($<$5\,ms per request) --- multipart parsing and HTTP proxying only.
  \item Madd analysis adds 200--600\,ms on top of the inference time when \texttt{reference\_text} is provided (one extra CTC forward pass for forced alignment).
\end{itemize}
\end{infobox}

\section{Processing Time Breakdown}

For a typical 10-second audio file with \texttt{reference\_text} provided (production CPU droplet, end-to-end through Caddy + Go + Python):

\begin{table}[H]
  \centering
  \caption{Request processing time breakdown (10\,s audio + madd, CPU)}
  \label{tab:perf_breakdown}
  \renewcommand{\arraystretch}{1.3}
  \begin{tabular}{lcc}
    \toprule
    \textbf{Stage} & \textbf{Duration (ms)} & \textbf{Percentage} \\
    \midrule
    Caddy TLS termination \& routing      & ${\sim}10$    & 0.1\,\% \\
    Go server parsing \& validation       & ${\sim}5$     & 0.1\,\% \\
    Network transfer (Go $\rightarrow$ Python) & ${\sim}15$ & 0.2\,\% \\
    Audio resampling (\texttt{librosa})   & ${\sim}45$    & 0.5\,\% \\
    RNNT inference (NeMo, FastConformer)  & ${\sim}8\,200$ & 90.0\,\% \\
    CTC forced alignment                  & ${\sim}500$   & 5.5\,\% \\
    Madd analyser                         & ${\sim}120$   & 1.3\,\% \\
    JSON serialisation \& return path     & ${\sim}30$    & 0.3\,\% \\
    \midrule
    \textbf{Total end-to-end}             & ${\sim}\mathbf{9\,100}$ & 100\,\% \\
    \bottomrule
  \end{tabular}
\end{table}

Model inference constitutes 90\,\% of total processing time, confirming that further performance work should target the ML backend (GPU upgrade, smaller checkpoint) rather than the API or analysis layers.

% ═══════════════════════════════════════════════════════════════════════════════
% CHAPTER 10: SCALABILITY AND LOAD HANDLING
% ═══════════════════════════════════════════════════════════════════════════════
\chapter{Scalability and Load Handling}

\section{Current Concurrency Model}

The Go API server handles concurrent HTTP connections natively through goroutines. The inference backend, however, intentionally serialises all model invocations via a module-level \texttt{asyncio.Semaphore(1)} plus a per-connection \texttt{asyncio.Lock} on the WebSocket path:

\begin{warnbox}[title=Concurrency Limitations]
\begin{itemize}
  \item \textbf{Global inference semaphore}: One inference at a time across the entire Python process. Prevents CPU over-subscription and non-deterministic output ordering. Concurrent requests are queued.
  \item \textbf{Single-worker Uvicorn}: The container runs one Uvicorn worker. Multi-process scaling would require multiple model copies --- ${\sim}500$\,MB of CPU RAM each --- which is infeasible on the baseline 8\,GB droplet.
  \item \textbf{Model memory footprint}: The FastConformer checkpoint occupies ${\sim}500$\,MB of CPU RAM (or ${\sim}1.2$\,GB of GPU VRAM in FP16). Loading multiple instances scales linearly.
\end{itemize}
\end{warnbox}

\section{Observed Behaviour Under Load}

Under concurrent load on the CPU production droplet the system exhibits a predictable serial-queue pattern:

\begin{table}[H]
  \centering
  \caption{Server behaviour under concurrent requests (CPU, 10\,s audio without madd)}
  \label{tab:scale_load}
  \renewcommand{\arraystretch}{1.3}
  \small
  \begin{tabular}{p{2.6cm}p{2.8cm}p{7.5cm}}
    \toprule
    \textbf{Concurrent Requests} & \textbf{Avg. Response Time} & \textbf{Behaviour} \\
    \midrule
    1   & ${\sim}8.5$\,s   & Normal \\
    3   & ${\sim}25$\,s    & Queued (serial execution) \\
    5   & ${\sim}42$\,s    & Queued, increased wait \\
    10+ & ${\sim}85$\,s+   & High risk of client-side timeout (5\,min ceiling at the Go layer) \\
    \bottomrule
  \end{tabular}
\end{table}

\section{Scalability Improvement Strategies}

\begin{goldbox}[title=Recommended Scalability Enhancements]
\begin{itemize}
  \item \textbf{GPU upgrade}: A single T4/L4 GPU droplet cuts final latency below 2\,s and lifts effective single-stream throughput by ${\sim}7\times$. The fastest unit-of-work improvement.
  \item \textbf{Horizontal scaling}: Run multiple droplets behind a load balancer. Each droplet is fully self-contained (Caddy issues its own cert), so adding capacity is one-shot.
  \item \textbf{Multi-worker Uvicorn (with care)}: Launch Uvicorn with \texttt{--workers N} only when RAM permits one model copy per worker. The semaphore stays per-process so concurrency is bounded by worker count.
  \item \textbf{Bounded request queue}: Add a token-bucket or in-memory queue at the Go layer to reject excess work with \texttt{503 Service Unavailable} rather than accumulating unbounded latency.
  \item \textbf{Batching on GPU}: On GPU hardware, batch multiple audio inputs in one forward pass to amortise the fixed kernel-launch overhead.
\end{itemize}
\end{goldbox}

% ═══════════════════════════════════════════════════════════════════════════════
% CHAPTER 11: SECURITY CONSIDERATIONS
% ═══════════════════════════════════════════════════════════════════════════════
\chapter{Security Considerations}

\section{Transport Layer}

Caddy terminates TLS using a Let's Encrypt certificate issued automatically on the first HTTPS request. Plain-text HTTP traffic on port 80 is redirected by Caddy's default behaviour. The Go and Python containers \emph{never} bind to a public address --- they expose ports only within the compose network.

\section{Input Validation}

The Go API gateway enforces multiple layers of input validation before any data reaches the inference backend:

\begin{emeraldbox}[title=Validation Rules]
\begin{itemize}
  \item \textbf{HTTP method}: \texttt{/api/v1/evaluate} accepts only \texttt{POST}; other methods return \texttt{405 Method Not Allowed}.
  \item \textbf{File size limit}: \texttt{r.ParseMultipartForm(50 << 20)} enforces a strict 50\,MB ceiling. Larger uploads return \texttt{400 Bad Request}.
  \item \textbf{File-format whitelist}: Only five audio extensions are accepted (\texttt{.wav}, \texttt{.mp3}, \texttt{.m4a}, \texttt{.flac}, \texttt{.ogg}). Unrecognised extensions are rejected before any read past the multipart header.
  \item \textbf{Limit-read of audio}: The handler wraps the form file in \texttt{io.LimitReader(file, 50<<20)} so a malformed multipart cannot stream past the declared ceiling.
  \item \textbf{Format check at the backend}: The Python \texttt{/evaluate} handler re-checks the extension on the server side as defence-in-depth.
\end{itemize}
\end{emeraldbox}

\section{Token Management}

Two distinct secrets are managed exclusively on the server side:

\begin{infobox}[title=Token Security Architecture]
\begin{itemize}
  \item \textbf{\texttt{HF\_TOKEN}}: Stored in \texttt{/opt/tahkik-server/.env} (mode 600) on the droplet. Injected into the Python container as an environment variable. Used to authenticate to HuggingFace when downloading private model weights. The Flutter app never sees it.
  \item \textbf{\texttt{HEALTH\_TOKEN}}: Generated on first provisioning (\texttt{head -c 32 /dev/urandom | base64}, 40 chars). Required as a Bearer token on every \texttt{/api/v1/system-health} call and as a \texttt{?token=} query parameter on the WebSocket variant. The Go middleware fails closed (\texttt{503}) if the env var is unset, so the endpoint cannot accidentally become public.
\end{itemize}
\end{infobox}

\section{Docker Socket Exposure}

The Go container mounts \texttt{/var/run/docker.sock} read-only (\texttt{:ro}) to enumerate sibling containers for the system-health endpoint. Read-only mode prevents container creation, image manipulation, or any state-changing API call. The distroless \texttt{nonroot} user (UID 65532) is added to the host's docker group (GID 999) at runtime via \texttt{group\_add} so the socket is readable without granting root inside the container.

\section{Abuse Prevention}

\begin{table}[H]
  \centering
  \caption{Current protections and proposed enhancements}
  \label{tab:security_measures}
  \renewcommand{\arraystretch}{1.3}
  \begin{tabular}{p{5cm}lp{5.5cm}}
    \toprule
    \textbf{Measure} & \textbf{Status} & \textbf{Description} \\
    \midrule
    TLS at edge (Let's Encrypt)            & Implemented & All client traffic encrypted by default \\
    File size cap (50\,MB)                 & Implemented & Prevents oversized-upload resource exhaustion \\
    Format whitelist                       & Implemented & Blocks arbitrary file uploads \\
    Inference semaphore                    & Implemented & Natural throttle: 1 concurrent inference \\
    HTTP timeout (5\,min)                  & Implemented & Prevents indefinite connection holding \\
    Bearer-token gate on system-health     & Implemented & Operator-only endpoint \\
    Distroless + nonroot Go container      & Implemented & No shell, no package manager, minimal attack surface \\
    Rate limiting                          & Proposed    & Token-bucket at Caddy or Go layer \\
    Per-client API keys                    & Proposed    & Identify and throttle individual mobile users \\
    \bottomrule
  \end{tabular}
\end{table}

%% --- continued in part 4 ---
% ═══════════════════════════════════════════════════════════════════════════════
% CHAPTER 12: REQUEST LIFECYCLE
% ═══════════════════════════════════════════════════════════════════════════════
\chapter{Request Lifecycle}

\section{End-to-End Request Flow}

The following describes the complete lifecycle of a single audio-evaluation request, from client submission to response delivery.

\phaseblock{Step 1 --- Client Submission}{The Flutter mobile application captures or selects an audio recording, optionally pairs it with the ayah text the user was asked to recite, and sends it as a \texttt{multipart/form-data} POST request to \texttt{https://<SITE\_HOST>/api/v1/evaluate}. The audio is attached under the \texttt{audio} form field; the reference (if any) under \texttt{reference\_text}.}

\phaseblock{Step 2 --- Caddy TLS \& Routing}{Caddy terminates the TLS handshake, runs the CORS preflight if applicable, and matches the request against the \texttt{@api} path matcher. The request is reverse-proxied to \texttt{go-server:8080} on the internal compose network. The CORS \texttt{Access-Control-Allow-Origin} response header is set from \texttt{\$DASHBOARD\_ORIGIN}.}

\phaseblock{Step 3 --- Go Server Validation}{The Go server parses the multipart form (enforcing the 50\,MB limit), extracts the \texttt{audio} field, validates the file extension against the whitelist, reads the binary data into memory through a limit-reader, and captures the optional \texttt{reference\_text}.}

\phaseblock{Step 4 --- Inference Forwarding}{The Go server constructs a fresh \texttt{multipart/form-data} request containing both \texttt{audio} and \texttt{reference\_text} (when supplied), targets \texttt{\$INFERENCE\_SERVER\_URL/evaluate}, and attaches the \texttt{Authorization: Bearer \$HF\_TOKEN} header if available. The HTTP client uses a 5-minute timeout.}

\phaseblock{Step 5 --- Backend Audio Processing}{The Python FastAPI handler receives the upload, writes it to a temp file under \texttt{/tmp}, and loads it via \texttt{librosa.load(sr=16000, mono=True)}. The audio array is clamped to the last 60 seconds.}

\phaseblock{Step 6 --- Model Inference}{The array is passed to \texttt{model.transcribe(\dots)} under the global inference semaphore. When \texttt{reference\_text} is present, the call requests timestamps and hypotheses so per-word and per-character timings can be extracted.}

\phaseblock{Step 7 --- Optional Madd Analysis}{The CTC head is run for forced alignment, the analyser computes per-rule errors, and the response payload is augmented with \texttt{errors[]} and \texttt{haraka\_ms}.}

\phaseblock{Step 8 --- Response Assembly}{The temp file is deleted in a \texttt{finally} block. The JSON result is returned to the Go server, which wraps it in the standard \texttt{\{"status":"success","data":\{\dots\}\}} envelope and forwards it back through Caddy to the Flutter client.}

\section{Concurrency Patterns}

Different components employ distinct concurrency strategies, each chosen for the workload it sees:

\begin{table}[H]
  \centering
  \caption{Concurrency mechanisms across components}
  \label{tab:concurrency}
  \renewcommand{\arraystretch}{1.3}
  \begin{tabular}{lll}
    \toprule
    \textbf{Component} & \textbf{Pattern} & \textbf{Rationale} \\
    \midrule
    Caddy           & Event loop (Go native) & High-throughput TLS reverse proxy \\
    Go API server   & Goroutines           & Lightweight; one per request \\
    Python REST     & Single Uvicorn worker & Backed by one model copy \\
    Python WS       & \texttt{asyncio} tasks & Bounded by per-connection lock \\
    NeMo inference  & \texttt{Semaphore(1)} & Single CPU/GPU shared resource \\
    \bottomrule
  \end{tabular}
\end{table}

% ═══════════════════════════════════════════════════════════════════════════════
% CHAPTER 13: FAILURE HANDLING
% ═══════════════════════════════════════════════════════════════════════════════
\chapter{Failure Handling}

\section{Error Categories and Responses}

The server distinguishes several categories of failure, each returning a specific HTTP status code and error message:

\begin{table}[H]
  \centering
  \caption{Failure modes and corresponding HTTP responses}
  \label{tab:failure_modes}
  \renewcommand{\arraystretch}{1.3}
  \begin{tabular}{p{4cm}cp{6cm}}
    \toprule
    \textbf{Failure Mode} & \textbf{HTTP Status} & \textbf{Response Action} \\
    \midrule
    Missing \texttt{audio} field      & 400 & \texttt{"audio field required"} \\
    Wrong HTTP method                 & 405 & \texttt{"method not allowed"} \\
    Oversized upload                  & 400 & \texttt{"file too large"} \\
    Unsupported audio format          & 400 & Reject with file extension info \\
    Inference timeout (5\,min)        & 504 & Return timeout error \\
    Backend (Python) returns 5xx      & 502 & Wrap and forward error \\
    Caddy cannot reach Python (cold) & 502 & Generic gateway error \\
    \texttt{HEALTH\_TOKEN} unset      & 503 & system-health endpoint fails closed \\
    Wrong Bearer token                & 401 & system-health rejects \\
    Internal Python exception         & 500 & Log + return generic message \\
    \bottomrule
  \end{tabular}
\end{table}

\section{Error Response Format}

All client-facing errors at the Go layer use a uniform envelope:

\begin{lstlisting}[caption={Uniform error response}]
{
  "status":  "error",
  "message": "<human-readable error description>"
}
\end{lstlisting}

This matches the success envelope's shape (\texttt{\{"status": "success", "data": \{\dots\}\}}) so clients can use a single discriminator field to branch.

\section{Resilience Mechanisms}

\begin{emeraldbox}[title=Implemented Resilience Patterns]
\begin{itemize}
  \item \textbf{Per-task exception isolation}: Errors in a WebSocket partial-window task are caught, logged with full traceback, and discarded --- they cannot tear down the WebSocket connection.
  \item \textbf{Temp-file cleanup}: The \texttt{/evaluate} handler deletes uploaded audio in a \texttt{finally} block regardless of whether inference succeeded or threw.
  \item \textbf{Per-container healthchecks}: Each container has a \texttt{HEALTHCHECK} directive (Go uses the \texttt{healthcheck} sub-command of its own binary; Python uses an inline urllib probe; Caddy's image healthcheck is sufficient).
  \item \textbf{\texttt{restart: unless-stopped}}: All three services restart automatically on crash unless explicitly stopped by an operator.
  \item \textbf{Cached system-health snapshot}: The 5-second cache absorbs traffic spikes from the dashboard's poll loop without thrashing \texttt{/proc} or Docker.
  \item \textbf{System-health WS keepalive}: The WebSocket sends \texttt{ping} frames and the server replies with \texttt{pong}; a missed pong tears down the connection and the dashboard reconnects.
  \item \textbf{Reset/stop session control}: The WebSocket protocol exposes explicit \texttt{reset} and \texttt{stop} messages; both cancel in-flight tasks and \texttt{reset} bumps a generation counter so stale partials are discarded silently.
\end{itemize}
\end{emeraldbox}

% ═══════════════════════════════════════════════════════════════════════════════
% CHAPTER 14: LOGGING AND MONITORING
% ═══════════════════════════════════════════════════════════════════════════════
\chapter{Logging and Monitoring}

\section{Logging Strategy}

The server employs a layered logging approach with structured log entries at each tier:

\begin{table}[H]
  \centering
  \caption{Logging implementation by component}
  \label{tab:logging}
  \renewcommand{\arraystretch}{1.3}
  \begin{tabular}{lp{5.5cm}p{5cm}}
    \toprule
    \textbf{Component} & \textbf{Log Output} & \textbf{Captured} \\
    \midrule
    Caddy           & Container stdout (Docker JSON)  & Access log + cert events \\
    Go server       & \texttt{log} package $\rightarrow$ stdout & Request lifecycle, errors, version banner \\
    Python REST     & \texttt{print(\dots, flush=True)}        & Model loading, request timing \\
    Python WS       & Same                                     & Connect/disconnect, partial/final emit, errors \\
    Madd analyser   & \texttt{print(\dots, flush=True)}        & Word alignment, haraka calibration, span counts \\
    \bottomrule
  \end{tabular}
\end{table}

All container logs are captured by Docker's \texttt{json-file} driver with 20\,MB rotation across 5 files (configured in \texttt{/etc/docker/daemon.json} by the provisioning script). Operators can tail any container with \texttt{docker compose logs -f <service>}.

\section{Metrics Exposure}

The Python container maintains lightweight in-process counters and rolling latency deques (200 samples each for partial and final paths). These are returned by \texttt{GET /health}:

\begin{lstlisting}[language=Python, caption={Metrics snapshot returned by Python /health}]
{
  "inference": {
    "requests_total":         123,
    "requests_failed":        2,
    "ws_connections_total":   45,
    "ws_connections_active":  3,
    "partial_p50_ms":         3450,
    "partial_p95_ms":         7100,
    "final_p50_ms":           8200,
    "final_p95_ms":           14500,
    "last_request_at":        "2026-...",
    "last_error_at":          null,
    "last_error":             null
  }
}
\end{lstlisting}

The Go server's \texttt{/api/v1/system-health} aggregator augments this with host metrics, container status, and version stamps for the operations dashboard.

\section{Observability Roadmap}

\begin{goldbox}[title=Proposed Observability Enhancements]
\begin{itemize}
  \item \textbf{Prometheus exposition}: Add a \texttt{/metrics} endpoint exporting the same in-process counters in Prometheus format for scraping by an external collector.
  \item \textbf{Structured JSON logs}: Replace the current line-based logs with single-line JSON entries so log shippers (Vector, Fluent Bit) can index them.
  \item \textbf{Request tracing}: Generate a unique trace ID at the Caddy layer and propagate through Go and Python logs to correlate per-request events.
  \item \textbf{DigitalOcean alert policies}: External backstop for the in-dashboard alerts --- four threshold policies (CPU, memory, disk, droplet-down) documented in \texttt{docs/do-alerts-setup.md}.
  \item \textbf{Centralized error reporting}: Forward Python tracebacks and Go panics to a Sentry or similar aggregator.
\end{itemize}
\end{goldbox}

\section{Dashboard Integration}

The Tahkik admin dashboard polls \texttt{/api/v1/system-health} every 5\,s for the snapshot view and opens a WebSocket to \texttt{/api/v1/system-health/ws} for the live metrics page. Both endpoints require the \texttt{HEALTH\_TOKEN}; the dashboard reads it from its own \texttt{VITE\_TAHKIK\_HEALTH\_TOKEN} environment variable and stores it only client-side (no propagation to third parties).

% ═══════════════════════════════════════════════════════════════════════════════
% CHAPTER 15: CONCLUSION
% ═══════════════════════════════════════════════════════════════════════════════
\chapter{Conclusion}

\section{Summary}

This report has documented the architecture, implementation, and operational concerns of the Tahkik inference server in its current NeMo-based form. The system delivers production-grade Arabic Quranic recitation transcription and Tajweed-rule analysis to the Flutter mobile application through a three-container Docker Compose stack: Caddy (TLS + routing), a Go API gateway (validation + system-health aggregation), and a Python FastAPI server hosting NVIDIA's FastConformer Hybrid CTC/RNNT model.

\section{Key Achievements}

\begin{infobox}[title=Project Highlights]
\begin{itemize}
  \item \textbf{Engine upgrade}: Successfully migrated from the legacy Whisper / CTranslate2 stack to NVIDIA NeMo FastConformer with no Flutter-side code change (only an \texttt{HF\_SPACE\_URL} swap).
  \item \textbf{Quranic-domain accuracy}: 6.55\,\% WER on the EveryAyah benchmark on the same canonical checkpoint that drives the inference path.
  \item \textbf{Real-time streaming}: \texttt{/ws/stream} delivers partial transcriptions on an 8-second sliding window with ${<}1$\,s p50 latency on a single T4 GPU.
  \item \textbf{Madd-rule grading}: The analyser implements eight Warsh/Azraq madd rules with auto-calibrated haraka length and CTC-derived per-letter timings.
  \item \textbf{One-command deploy}: \texttt{provision-droplet.sh} stands up a fresh production droplet end-to-end (Docker + TLS + model cache) in 10--15 minutes.
  \item \textbf{Operator dashboard}: \texttt{/api/v1/system-health} surfaces host, container, and inference metrics under a Bearer-token-gated endpoint with a WebSocket variant for live updates.
\end{itemize}
\end{infobox}

\section{Future Enhancements}

\begin{goldbox}[title=Roadmap]
\begin{itemize}
  \item \textbf{Performance}
    \begin{itemize}
      \item Migrate to GPU droplets (T4/L4) as traffic warrants.
      \item Quantise the NeMo checkpoint to INT8 for CPU inference (requires NVIDIA's quantisation tooling).
      \item Batch concurrent inferences on GPU.
    \end{itemize}
  \item \textbf{Reliability}
    \begin{itemize}
      \item Add rate limiting at the Caddy layer.
      \item Add a bounded request queue at the Go layer with explicit 503 responses.
      \item Wire Sentry or equivalent for error aggregation.
    \end{itemize}
  \item \textbf{Tajweed analysis}
    \begin{itemize}
      \item Extend rules beyond madd: idgham, ikhfa, qalqala, ghunna detection.
      \item Add tafkhim/tarqiq detection for the lam and ra letters.
      \item Multi-wajh grading mode (accept 2 / 4 / 6 counts) for Badal, Aarid, Lin instead of locked qasr.
    \end{itemize}
  \item \textbf{Operations}
    \begin{itemize}
      \item Multi-region deployments behind a global load balancer.
      \item Automated blue/green deploys via the GitHub Actions workflow.
      \item Prometheus + Grafana stack for long-term metrics retention.
    \end{itemize}
\end{itemize}
\end{goldbox}

\section{Final Remarks}

\begin{center}
\begin{tikzpicture}
  \node[
    draw=teal!50!black,
    fill=tealbg,
    rounded corners=10pt,
    inner sep=15pt,
    text width=14cm,
    align=center,
    text=teal!60!black
  ] {
    \large\bfseries Production-Ready Inference Stack \\[6pt]
    \normalsize
    The Tahkik server stack runs reliably on a single-host Docker Compose deployment, \\
    with a production CPU baseline at \texttt{s-4vcpu-8gb} and a GPU-overlay one-liner away. \\
    The architecture cleanly separates TLS, REST routing, and ML inference, \\
    making each layer independently upgradable.
  };
\end{tikzpicture}
\end{center}

% ═══════════════════════════════════════════════════════════════════════════════
% REFERENCES
% ═══════════════════════════════════════════════════════════════════════════════
\chapter*{References}
\addcontentsline{toc}{chapter}{References}

\begin{thebibliography}{99}

\bibitem{nemo}
NVIDIA Corporation. \emph{NeMo: A Toolkit for Conversational AI}. \\
\url{https://github.com/NVIDIA/NeMo}

\bibitem{fastconformer}
Rekesh, D., Koluguri, N. R., \emph{et al.} (2023). \emph{Fast Conformer with Linearly Scalable Attention for Efficient Speech Recognition}. ASRU 2023. \\
\url{https://arxiv.org/abs/2305.05084}

\bibitem{nemoarabic}
NVIDIA. \emph{stt\_ar\_fastconformer\_hybrid\_large\_pcd\_v1.0 model card}. \\
\url{https://huggingface.co/nvidia/stt_ar_fastconformer_hybrid_large_pcd_v1.0}

\bibitem{rnnt}
Graves, A. (2012). \emph{Sequence Transduction with Recurrent Neural Networks}. ICML 2012 Representation Learning Workshop. \\
\url{https://arxiv.org/abs/1211.3711}

\bibitem{ctc}
Graves, A., Fern\'andez, S., Gomez, F., Schmidhuber, J. (2006). \emph{Connectionist Temporal Classification: Labelling Unsegmented Sequence Data with Recurrent Neural Networks}. ICML 2006.

\bibitem{torchaudio_align}
PyTorch Team. \emph{Forced Alignment with Wav2Vec2 / CTC --- \texttt{torchaudio.functional.forced\_align}}. \\
\url{https://pytorch.org/audio/stable/generated/torchaudio.functional.forced_align.html}

\bibitem{fastapi}
Ram\'irez, S. \emph{FastAPI: Modern, fast (high-performance) web framework}. \\
\url{https://fastapi.tiangolo.com}

\bibitem{uvicorn}
Christie, T. \emph{Uvicorn --- The lightning-fast ASGI server}. \\
\url{https://www.uvicorn.org}

\bibitem{caddy}
Matt Holt et al. \emph{Caddy --- The Ultimate Server with Automatic HTTPS}. \\
\url{https://caddyserver.com}

\bibitem{gohttp}
The Go Authors. \emph{Package net/http}. \\
\url{https://pkg.go.dev/net/http}

\bibitem{coderws}
Coder. \emph{github.com/coder/websocket --- minimal and idiomatic WebSocket library for Go}. \\
\url{https://github.com/coder/websocket}

\bibitem{docker}
Docker, Inc. \emph{Docker Engine documentation}. \\
\url{https://docs.docker.com/engine/}

\bibitem{compose}
Docker, Inc. \emph{Docker Compose specification}. \\
\url{https://docs.docker.com/compose/}

\bibitem{librosa}
McFee, B. \emph{librosa: Audio and music signal analysis in Python}. \\
\url{https://librosa.org}

\bibitem{warsh}
\emph{Tariq al-Azraq (Warsh `an Nafi' transmission)} --- riwaya documentation in the Tajweed literature. Standard reference works on Madinah Mushaf and Tariq Azraq variants.

\bibitem{tahkikserver}
Ben Hadj Mohamed. \emph{Tahkik Server Source Code Repository}. \\
\url{https://github.com/mohamedbenhadjer/tahkik-server}

\bibitem{everyayah}
EveryAyah. \emph{Quranic Audio Dataset}. \\
\url{https://everyayah.com}

\end{thebibliography}

\vfill
\begin{center}
\textcolor{slategray}{\itshape End of Report}
\end{center}

